\section{Timing, efficiency and error handling}
The scheduler has idle, discovery, convert-start, wait, read and completed phases.
Each call advances one small step. Discovery performs one ROM search, reading
performs one addressed scratchpad, and waiting returns immediately until the
deadline. Broadcast conversion begins all externally powered sensors together.
The 760 ms deadline includes command execution and millisecond rounding around
the sensor's 750 ms maximum. Deadlines use signed subtraction and handle wrap;
service gaps must remain below half the 32-bit millisecond range.

\subsection{Real execution costs}
Nonblocking conversion does not make the 1-Wire transport asynchronous. At
standard speed, a ROM search uses approximately 15 ms and an addressed read
approximately 12 ms of foreground time. A complete array cycle grows with sensor
count. Interrupts can execute during recovery and the long conversion wait.
Critical requested delay durations are about 70 us for reset sampling, 60 us for
write zero, 6 us for write one and 13 us for read sampling. Callback and register
overhead add to these figures; they are not measured upper bounds on this board.
For stricter real-time requirements replace the transport with a hardware master
or timer/DMA design and validate the actual waveform with an oscilloscope.

\subsection{Validate before consuming}
ROM and scratchpad CRC8 use the Dallas polynomial in reflected form 0x8C.
Scratchpad configuration and fixed reserved bytes are checked to reject an
all-zero line response, which would otherwise have a mathematically valid zero
CRC. Signed raw decoding masks undefined bits for 9--11 bit resolution. Converted
values outside -55 to 125 C are rejected. At conversion start, every array output
becomes NAN and BUSY; success changes its status to OK. Missing, timeout, CRC and
range failures retain NAN. Never compare NAN with equality; inspect status or
use the standard \texttt{isnan} function.

\subsection{Common failure diagnoses}
No device: check power, ground, pull-up and PB10. CRC error: inspect supply noise,
cable loading and timing; a disconnected addressed sensor can also cause CRC
error when another sensor still answers reset. Timeout: check sensor power,
conversion readiness and a line held low. Invalid argument: check the ROM CRC,
buffer lifetime, count and duplicate addresses. Temperature 85 C: verify that a
conversion completed and investigate supply resets, while accepting that 85 C
may be real. A low idle line is rejected before reset; a line that remains low
after the presence window is also rejected.

A configurable idle gap follows each cycle; \texttt{sample\_interval\_ms}
starts at 250 ms and can be adjusted before servicing. This also limits retries
when no sensor is attached. The conversion deadline is armed on the next call
after command transmission, so an interrupt during the command cannot shorten it.
